\section{Where the fixed cost lives: the kernel boundary, the shards
and the memory roof}
\label{sec:iopath}

\begin{figure*}[t]
\centering
\lrfigure{iopath}{\linewidth}
\caption{Per-packet cost and measured system calls per packet against batch
depth, at payload \protoMaxPayload\ bytes. Syscall counts are counted inside
each backend rather than derived from the design, because a batch that returns
short does not issue the number of calls the design intends.}
\label{fig:iopath}
\end{figure*}

Four transports carry the same hot path. \texttt{posix} issues one receive and
one send per packet. \texttt{mmsg} issues one of each per batch.
\texttt{uring} batches submissions through a shared ring \cite{iouring}.
\texttt{uring-sqpoll} adds a kernel-side submission poller, so submission costs
no system call at all. That is as close to a kernel-bypass data path
\cite{netmap,barbette2018} as an unprivileged process reaches without a
userspace driver.

The last of those is the point of the experiment. Without it, the system-call
share of $a$ can only be extrapolated from the shape of the batching curve. With
it, the quantity is measured: at batch depth one, \texttt{posix} costs
\ioPosixCycles\ cycles per packet and \texttt{uring-sqpoll} costs
\ioSqpollCycles, a difference of \ioBoundaryCycles\ cycles.

\noindent\textbf{When that difference is not a bound.}
A submission poller is a second runnable thread. On a host with fewer usable
cores than the poller needs, it competes with the data plane for the same core,
and the calls it removes cost more in scheduling than they saved. The harness
detects and records that condition rather than inferring it, because a negative
saving is a property of the host and not of the mechanism, and presenting it as
either a bug or a benefit would be wrong. The dataset records the condition as
follows: \ioSqpollNote. The measured difference is nevertheless negative, so on
this run the polled path bounds nothing. What it measures is the poller
competing for a machine with \envCores\ physical cores, and
Table~\ref{tab:iopath} shows \texttt{uring-sqpoll} approaching the other
backends only at the deepest batch in the grid.

\noindent\textbf{Batching.}
Table~\ref{tab:iopath} gives the full grid. The batched backends converge as the
depth grows, which is expected: past the point where the amortised system call
is small next to the cipher, the remaining differences are in how each ring
reports completions. On this host the best batched configuration was
\ioBestBatched\ at \ioBestBatchedCycles\ cycles per packet.

The measured syscall counts are worth reading alongside the costs. They confirm
that each backend does what its design says: two calls per packet for
\texttt{posix}, $2/B$ for \texttt{mmsg}, and none on the submission side for
\texttt{uring-sqpoll}. That is the check that the configuration under
measurement is the configuration intended.

\begin{table}[t]
\caption{Transport grid at payload \protoMaxPayload\ bytes.}
\label{tab:iopath}
\centering\footnotesize
\input{generated/table_iopath}
\end{table}

\begin{figure}[t]
\centering
\lrfigure{scaling}{\columnwidth}
\caption{Strong speedup and weak efficiency against the number of shards. Shards
share nothing mutable, so what this measures is contention for the memory system
rather than for a lock.}
\label{fig:scaling}
\end{figure}

\noindent\textbf{Widening the data plane.}
Each shard owns its socket, key schedule, replay window and arena, and shares
nothing mutable with any other. The scaling question is therefore not whether a
lock contends, but what the machine does when $n$ identical independent
pipelines run at once: last-level cache, memory bandwidth, and on a
hyper-threaded host the sharing of one physical core's execution ports.

Two laws are reported, because they answer different questions and quoting one
for the other is a standard way to overstate a result. \emph{Strong} scaling
fixes the total packet count and divides it, and answers whether adding a core
helps with a given amount of traffic. \emph{Weak} scaling gives each shard the
same count whatever $n$ is, and answers whether the machine can absorb $n$ times
the traffic with $n$ times the cores.

The worker pool is expressed with OpenMP where the compiler provides it and with
\texttt{std::thread} otherwise; on this run it was \scalePool. Threads are
created and shards constructed and warmed before the timed region opens, and all
workers wait at a barrier, so that thread creation is not charged to the data
plane and no worker measures an uncontended machine because it started first.

Fitting Amdahl's law to the strong-scaling curve gives a serial fraction of
\scaleSerialFraction\ over \scaleMaxShards\ shards. Table~\ref{tab:scaling}
gives the points. The cross-node half of this experiment, which runs one MPI
rank per node to see whether the per-packet cost stays put when the ranks are on
separate hosts, did not run here: its fragment is the one missing from the
dataset, and Section~\ref{sec:threats} says what that costs the argument.

\begin{table}[t]
\caption{Scaling of the shard pool.}
\label{tab:scaling}
\centering\footnotesize
\input{generated/table_scaling}
\end{table}

\begin{figure}[t]
\centering
\lrfigure{memroof}{\columnwidth}
\caption{Triad bandwidth against working-set size. The plateaus are the cache
hierarchy; the rightmost is the sustained figure the roof is computed from.}
\label{fig:memroof}
\end{figure}

\noindent\textbf{The memory roof.}
A data plane cannot move payload bytes faster than the memory system supplies
them, divided by the number of times each byte is touched. A STREAM triad
\cite{mccalpin1995} over a geometric range of working sets measures the
numerator directly, so the scaling curve above can be read against a roof
\cite{williams2009} rather than against nothing.

The largest working set measured sustains \memSustainedGBps\ GB/s. Each payload
byte makes \memPasses\ passes through memory in the hot path: the kernel copy in
on receive, the read and write of the decryption, the read and write of the
re-encryption, and the kernel copy out on send, each counted as one read and one
write. Dividing gives a roof of \memRoofGBps\ GB/s of payload, the rate this
data plane cannot exceed on this machine however fast the cipher becomes.
Figure~\ref{fig:memroof} shows the plateaus of the cache hierarchy that produce
that number, and therefore which level a packet buffer of a given size is served
from.
